\documentclass[11pt]{article}
\usepackage[margin=1in]{geometry}
\usepackage{amsmath,amssymb,amsthm,hyperref}
\newtheorem{theorem}{Theorem}
\newtheorem{proposition}[theorem]{Proposition}
\newcommand{\ii}{\mathrm i}
\newcommand{\Span}{\operatorname{span}}
\newcommand{\perm}{\operatorname{perm}}
\title{A sharp Fourier-family bound for the averaging lemma}
\author{Research note}
\date{7 October 2026}
\begin{document}
\maketitle
\noindent\textbf{Status.} The desired implication for \emph{all} order-six
Hadamards, $d_F(H)\ge 1/30\Rightarrow M(H)\ge5$, remains unproved here.
The follow-up \texttt{PRACTICAL\_LOCAL.tex} now proves the actual cutoff-$5$
implication whenever an equivalent representative has horizontal phase
distance at most $3/10$ from $F_6$, with an exact-arithmetic certificate.
Global confinement into that neighbourhood remains open. The extremely
small-cutoff calculations later in this older note do not meet the practical
target and are retained only as previous research.
This note proves its sharp version for the entire Fourier family and its
transpose, proves that the conjectured global constant is a rigorous upper
bound, replaces numerical tangent minimization by exact formulas, and gives
a robust rational replacement for the final irrationality argument in MC1.
The follow-up below proves global confinement at an explicit small cutoff
using three-row strips, without importing global MC1. That cutoff remains
far too small for the proposed practical bound of 5.
The orbit-distance argument and trigonometric estimate are analytic;
the finite corner identities are checked by exact arithmetic.

All matrices below are raw Hadamards: their entries have modulus one and
$HH^*=6I$. Projectors use the normalized columns $H_j/\sqrt6$.
Let $M(H)=\max_c R_c(H)$, over the 400 positional $3\times3$ corners,
with the invariant and the monomial-orbit distance $d_F$ defined in the handoff.

\section{The Fourier-family theorem}
Let $F=(e^{2\pi\ii jk/3})_{j,k=0}^2$, $Z=\operatorname{diag}(1,e^{\ii\alpha},e^{\ii\beta})$, and
\[
 N(\alpha,\beta)=\begin{pmatrix}F&F\\FZ&-FZ\end{pmatrix}.
\]
We use the repository's exact corner census:
\begin{align}
 X&=2\beta-\alpha,&Y&=\alpha+\beta,&T&=X-Y,\notag\\
 M(N)&=36\bigl(|\sin X|+|\sin Y|+|\sin T|\bigr)^2.\tag{1}
\end{align}
The complete multiset is zero 256 times and
$36(\sin X+s\sin Y+t\sin T)^2$, for $s,t\in\{-1,1\}$,
each 36 times. If $X,Y\in\pi\mathbb Z$, the matrix is equivalent to $F_6$.
These statements, including explicit equivalences at the twelve parameter
points modulo $2\pi$, are independently replayed by \texttt{verify.py}.

\begin{theorem}
For a Fourier-family matrix, its transpose, or any matrix equivalent to
either, if $M\le36$, then
\[
 \boxed{\quad d_F\le
 \sin\!\left(\frac16\arcsin\frac{\sqrt M}{12}\right).\quad} \tag{2}
\]
Consequently, throughout these families,
\[
 d_F\ge\frac1{30}\quad\Longrightarrow\quad
 M\ge c_*:=144\sin^2\!\left(6\arcsin\frac1{30}\right)
 =\frac{81966334263296}{14416259765625}>5.
 \tag{3}
\]
The constant $c_*=5.685686550872\ldots$ is attained.
\end{theorem}

\begin{proof}[Proof of the bound]
Put $S=\sqrt M/6\le1$. Choose integers $m,n$ such that
$x=X-m\pi$ and $y=Y-n\pi$ belong to $[-\pi/2,\pi/2]$.
Since $|\sin x|+|\sin y|\le1$, we have
$\sin^2|x|+\sin^2|y|\le1$, hence $|x|+|y|\le\pi/2$.
Therefore all three distances between $0,x,y$ are at most $\pi/2$.

Let $t$ be the largest of these distances. The other two are $u$ and
$t-u$ for some $0\le u\le t$. Consequently
\[
 S=\sin t+\sin u+\sin(t-u)\ge 2\sin t,
 \qquad t\le\arcsin(S/2). \tag{4}
\]
Here $\sin u+\sin(t-u)\ge\sin t$ follows directly from concavity of
$\sin$ on $[0,\pi/2]$, or from its addition formula.

The reference point
$\alpha_0=(2n-m)\pi/3$, $\beta_0=(m+n)\pi/3$ is equivalent to $F_6$.
Write
\[
 a=\alpha-\alpha_0=\frac{2y-x}{3},\qquad
 b=\beta-\beta_0=\frac{x+y}{3}.
\]
The span $r=\max(0,a,b)-\min(0,a,b)$ satisfies $r\le2t/3$:
each of $3a=y+(y-x)$, $3b=x+y$, and $3(a-b)=(y-x)-x$
has absolute value at most $2t$, using $|x|,|y|,|x-y|\le t$.

Let $\lambda$ be the midpoint of the interval spanned by $0,a,b$.
Multiply the lower three rows by $e^{-\ii\lambda}$ and the two columns
of phase class $d_j\in\{0,a,b\}$ by $e^{-\ii(d_j-\lambda)/2}$.
Relative to $N(\alpha_0,\beta_0)$ every entry now has phase error
$\pm(d_j-\lambda)/2$, of absolute value at most $r/4$.
Every normalized column overlap therefore has real part at least
$\cos(r/4)$. For unit vectors, projector distance is
$\sqrt{1-|\langle v,w\rangle|^2}$; thus
\[
 d_F\le\sin(r/4)\le\sin(t/6)
 \le\sin\!\left(\frac16\arcsin(S/2)\right).
\]
The same entrywise bound controls every normalized row, so it also proves
the assertion after transposition. No transpose invariance of $d_F$ is
assumed. Monomial equivalence preserves $d_F$ and $M$.

For (3), if $M>36$ there is nothing to prove. Otherwise invert (2),
using monotonicity on the displayed angle ranges.
\end{proof}

\section{Exact attainment, including minimization over the orbit}
Index rows and columns from zero. Let $t_1$ be one on rows $1,3,5$
and columns $1,4$, and zero elsewhere; let $t_2$ use columns $2,5$
instead. Define
\[
 H(a)=F_6\odot\exp\bigl(\ii a(t_1-t_2)\bigr).
\]
This is an affine Fourier-family Hadamard: same-parity row pairs retain
their Fourier orthogonality, and for opposite-parity row pairs the terms
in columns $j,j+3$ cancel. The exact corner census gives
\[
 M(H(a))=144\sin^2(3a),\tag{5}
\]
with this value appearing 72 times and zero appearing 328 times.

\begin{proposition}
For $0\le a\le1/12$,
\[
 d_F(H(a))=\sin(a/2).\tag{6}
\]
In particular the example at $a=2\arcsin(1/30)$ proves both attainment
in (3) and the rigorous global upper bound
$\inf_{d_F(H)\ge1/30}M(H)\le c_*$.
\end{proposition}

\begin{proof}
Matching the columns to $F_6$ without row phases gives overlaps
$1$ or $\cos(a/2)$ in absolute value, hence the upper bound.

First fix a matching whose dephased reference is exactly $F_6$, and
allow arbitrary row phases. Let $A$ and $B$ be the means of their
phase factors over the even and odd rows, respectively; $|A|,|B|\le1$.
The two extreme column classes have squared overlaps
$|A+Be^{\ii a}|^2/4$ and $|A+Be^{-\ii a}|^2/4$.
Their mean is
\[
 \frac{|A|^2+|B|^2+2\operatorname{Re}(\overline A B)\cos a}{4}
 \le\frac{1+\cos a}{2}=\cos^2(a/2).
\]
At least one column thus has projector distance at least $\sin(a/2)$.

It remains to rule out a better permutation of the Fourier reference.
Suppose such a matching has distance $\epsilon<\sin(a/2)\le a/2$.
Choose column phases aligning the two normalized columns. Their
Euclidean distance is at most $\sqrt2\epsilon$, so each raw entry
differs by at most $\sqrt{12}\epsilon$. Two unit complex numbers with
chord distance $q$ have circular phase distance at most $\pi q/2$.
Thus the entrywise phase discrepancies are at most
$\pi\sqrt3\epsilon$; after dephasing they are at most
$4\pi\sqrt3\epsilon$.

All entries of a dephased matrix in the $F_6$ orbit are sixth roots
of unity. Since $H(a)$ is itself dephased and differs in phase from
$F_6$ by at most $a$, our hypothetical reference differs from $F_6$
in circular phase by less than
\[
 a+4\pi\sqrt3\epsilon\le(1+2\pi\sqrt3)a<1<\pi/3.
\]
For the strict inequality use $a\le1/12$, $\pi<22/7$, and
$\sqrt3<7/4$. Distinct sixth roots are separated by at least $\pi/3$,
so this reference must dephase to $F_6$. The preceding lower bound
rules out its being better. The case $a=0$ is immediate.

Finally, $\sin(1/24)\ge1/24-(1/24)^3/6>1/30$, so the required
$a=2\arcsin(1/30)$ is less than $1/12$.
\end{proof}

\section{The sharp tangent inequality requires no manifold assumption}
The linearized dephased Hadamard equations at $F_6$ have kernel
spanned by $t_1,t_2,t_1^T,t_2^T$. For a tangent
$v=a t_1+b t_2+c t_1^T+d t_2^T$, put
\[
 Q(a,b)=a^2-ab+b^2,\qquad
 A(a,b)=\max\{|a+b|,|2a-b|,|a-2b|\},
\]
and let $r(a,b)$ be the span of $0,a,b$.
Write $M_2(v)$ for the maximum squared linear signed-corner invariant,
and $d_1(v)$ for the first-order projector distance, allowing row phases.

\begin{proposition}
Exactly,
\begin{align}
 M_2(v)&=144\bigl(A(a,b)+A(c,d)\bigr)^2,\tag{7}\\
 d_1(v)^2&=r(a,b)^2/16+Q(c,d)/18,\tag{8}\\
 \|v_\perp\|^2&=2\bigl(Q(a,b)+Q(c,d)\bigr).\tag{9}
\end{align}
Consequently
\[
 M_2(v)\ge216\|v_\perp\|^2,\qquad
 M_2(v)\ge5184d_1(v)^2.\tag{10}
\]
Both constants are sharp, for example at $(a,b,c,d)=(1,-1,0,0)$.
\end{proposition}

\begin{proof}
For a phase block $B$, set
$\mathcal D_B=2\operatorname{Im}(\perm(B)\det(B)^2/\prod B_{ij})$.
On a Hadamard, $M=\max_B\mathcal D_B^2$ by the complementary-block
identity. Exact differentiation at $F_6$ gives 220 zero forms,
the forms $12\ell(a,b)$ and $12\ell(c,d)$ each twelve times, and
all 36 forms $12(\ell(a,b)+\ell'(c,d))$ once, where
\[
 \ell,\ell'\in\{\pm(a+b),\ \pm(2a-b),\ \pm(a-2b)\}
\]
with the appropriate argument pair substituted. This proves (7).

For (8), remove row and column means. Let
$h=(-1,1,-1,1,-1,1)^T/2$, and let $u,w$ be the mean-zero versions
of $(0,a,b,0,a,b)^T$ and $(0,c,d,0,c,d)^T$. Then
\[
 v_\perp=hu^T+wh^T,\quad h\perp w,\quad
 \|h\|^2=3/2,\quad \|w\|^2=4Q(c,d)/3.
\]
Optimizing row phases to first order is precisely translating the six
column vectors of this matrix to minimize their maximum squared norm,
divided by six. These vectors are the three pairs $u_jh\pm w/2$.
The four extreme points form a rectangle. Its smallest enclosing ball
has radius squared $\|h\|^2r(a,b)^2/4+\|w\|^2/4$.
This proves (8); (9) follows by taking the Frobenius norm.

If three real numbers sum to zero and their maximum absolute value is
$m$, their sum of squares is at most $2m^2$.
Applying this to the mean-zero versions of $0,a,b$ gives
$A(a,b)^2\ge3Q(a,b)$. Also $A(a,b)\ge3r(a,b)/2$.
Thus
\begin{align*}
 144(A(a,b)+A(c,d))^2
 &\ge432(Q(a,b)+Q(c,d)),\\
 144(A(a,b)+A(c,d))^2
 &\ge324r(a,b)^2+432Q(c,d)\\
 &\ge324r(a,b)^2+288Q(c,d)=5184d_1(v)^2.
\end{align*}
\end{proof}

These are tangent statements, not estimates on a specified finite
neighbourhood. They imply the sharp asymptotic lower bound
$\liminf_{H\to F_6,\,H\not\sim F_6} M(H)/d_F(H)^2\ge5184$.
Indeed, choose 21 independent Hadamard equations at $F_6$.
Their implicit-function graph contains every sufficiently close actual
Hadamard, with normal displacement quadratic in its tangent displacement.
Taylor expansion gives $M=M_2+O(\|v\|^3)$ and
$d_F^2=d_1^2+O(\|v\|^3)$; the orbit-separation argument excludes
different discrete matchings sufficiently near $F_6$.
This reasoning does not assert that the remaining nine equations vanish
throughout the graph, or that the Hadamard set is a smooth four-manifold.

\section{A robust simplification of MC1's counting argument}
For any fixed three rows of a Hadamard, complementary column triples have
the same $p$: their three off-diagonal row overlaps are negatives of one
another. Thus the twenty corners form ten equal-$p$ pairs. The exact sum
rule $\sum_{20}p=108$ gives
\[
 \sum_{j=1}^{10}p_j=54 \tag{11}
\]
for one representative of each pair.

\begin{proposition}
Let $\kappa>0$ solve $3\kappa^3+\kappa^2+\kappa-1=0$, and put
$p_*=9-4/(1+\kappa^2)$. In every three-row strip of a Hadamard, some
corner satisfies
\[
 \operatorname{dist}\bigl(p,\{0,5,p_*\}\bigr)>1/32.\tag{12}
\]
In particular, the final counting contradiction in MC1 does not require
irrationality of $p_*$.
\end{proposition}
\begin{proof}
The cubic is strictly increasing on $[0,\infty)$ and has opposite signs
at $469/1000$ and $47/100$. Substitution in $9-4/(1+\kappa^2)$ gives
\[
 5.72<p_*<5.73.\tag{13}
\]
Suppose all ten representatives in (11) were within $1/32$ of
$\{0,5,p_*\}$. Choose corresponding centers $q_j$. Then
$|\sum q_j-54|\le10/32=5/16$.
If any center is zero, $\sum q_j\le9(5.73)=51.57$, a contradiction.
Otherwise, if $n$ centers equal $p_*$, their sum is $50+n(p_*-5)$.
For $n\le5$ this is less than $53.65$; for $n\ge6$ it is greater than
$54.32$. Both gaps from 54 exceed $5/16$, again a contradiction.
All decimals in this proof are terminating rational numbers.
\end{proof}

The sum rule itself is elementary. For a row pair, write the six unit
phase ratios as $z_j$, so $\sum z_j=0$. Expanding
$\sum_{|T|=3}|\sum_{j\in T}z_j|^2$ gives
$6\binom52+4\sum_{j\ne k}z_j\overline z_k=60-24=36$.
Adding the three row pairs gives 108.

This proposition isolates the quantitative task in MC1. Its exact
classification and routing steps say that an all-critical, non-Fourier
carrier would have every corner in the three scalar alternatives.
For a practical cutoff, those steps must be stabilized: it would suffice
to show that a carrier with $M<5$ outside the $1/30$ tube has a row strip
whose every $p$ is within $1/32$ of the alternatives. The proposition
would then give an immediate contradiction. No such stability implication
is proved by the exact-zero classification alone.

\section{What remains for the global lemma}
MC1 and compactness already prove the lemma with some unspecified
$c_0>0$: the Hadamard set is compact, $d_F$ is continuous, and $M$ has
zero set exactly the $F_6$ orbit. Its minimum on the closed set
$\{d_F\ge1/30\}$ is therefore strictly positive.
Relative to MC1 and the published effective-minimum theorem, the existing
\texttt{ProvenResults/EffectiveConfinement/PROOF.md} supplies the explicit
but impractical choice
\[
 c_0=2^{-N},\qquad N=2\cdot181\cdot26\cdot76^{26}.
\]
This is an existing result, not an improvement obtained in this note.

Neither compactness, the Fourier-family theorem, nor the tangent bound
proves $c_0=5$ for arbitrary Hadamards. The missing assertion is the
global implication
\[
 \boxed{\quad M(H)<5\quad\Longrightarrow\quad d_F(H)<1/30\quad}
\]
away from the explicit Fourier families. A Taylor estimate at $F_6$
does not exclude a separate region with small $M$. A quantitative
extension of the reconstruction's critical-strip argument, or a certified
cover of that region, is still needed.

\section{An explicit stability theorem from three-row strips}
Here ``small $R$'' means small \emph{maximum} $M(H)$, not just a small
value at one chosen corner. Set
\[
 N_0=2\cdot866\cdot11\cdot80^{11}
 =16365543384678400000000000,\qquad m_0=2^{-N_0}.
\]
\begin{theorem}[Explicit small-discriminant confinement]
Every order-six Hadamard satisfies
\[
 0\le M(H)<m_0\quad\Longrightarrow\quad d_F(H)<1/30.\tag{14}
\]
This proves the requested proximity-or-approximate-alternatives implication
for a specified sufficiently small cutoff, by forcing its proximity branch.
It does \emph{not} replace $m_0$ by 5.
\end{theorem}
The proof below uses the critical-strip classification and its two endpoint
exclusions, the near-$F_3$ stability estimate, and the effective-minimum
theorem of Jeronimo--Perrucci--Tsigaridas. Global MC1 is not an input.

\subsection*{Step A: a Fourier-family neighbourhood at the cutoff 5}
The signed corner polynomial $D_B=\ii\mathcal D_B$ has 48 Laurent terms
with coefficients $\pm1$. In the nine raw entries, 36 exponent vectors
have $\ell^1$ norm 8 and twelve have norm 12. Telescoping products of
unit numbers therefore gives, whenever $\max_{ij}|H_{ij}-K_{ij}|\le e$,
\[
 |\sqrt{M(H)}-\sqrt{M(K)}|\le432e.\tag{15}
\]
The estimate is valid in a common monomial frame; the matrices need not
be dephased. Normalized column distance, and hence the corresponding
projector distance, is at most $e$.

Suppose $M(H)\le5$, $K$ belongs to a Fourier family or its transpose,
and $e\le1/3600$. Then $\sqrt{M(K)}<56/25+432/3600=59/25$.
For $0\le x\le1/5$,
\[
 \arcsin x\le x+x^3/5.
\]
Indeed $(1+3t^2/5)^2(1-t^2)-1
=t^2(5-21t^2-9t^4)/25\ge0$ on this interval; integrate the resulting
bound on $(1-t^2)^{-1/2}$. Combining (2), (15) and the projector triangle
inequality, with $z=59/300$, yields
\[
 d_F(H)\le\frac1{3600}+\frac{z+z^3/5}{6}
 =\frac{26980379}{810000000}<\frac1{30}.\tag{16}
\]
This is a practical neighbourhood theorem at $M\le5$, independent of
the much smaller cutoff in (14).

We also use the proved near-$F_3$ estimate from
\texttt{ProvenResults/FourierNeighbourhood/SeedLifting/NEAR\_F6\_CLOSURE.md},
Section 3: a dephased $F_3$ corner with maximum entry error $h\le2^{-42}$
places the whole matrix within $1800\sqrt h$ entry error of a Fourier
family, its transpose, or Tao. The estimate controls both orientations.
For $h\le2^{-46}$ its error is at most $1800/2^{23}<1/3600$.
Tao is impossible when $M(H)\le5$, since (15) would give
$\sqrt{243}\le\sqrt5+432/3600<3$.
Thus
\[
 \left.\begin{array}{c}M(H)\le5\\
 \text{some dephased corner is within }2^{-46}\text{ of }F_3
 \end{array}\right\}\quad\Longrightarrow\quad d_F(H)<1/30.
 \tag{17}
\]

\subsection*{Step B: quantitative routing from a nearly singular corner}
We first prove a phase-block estimate. For $0<h\le1$, any $3\times3$
phase matrix $E$ with $|\det E|\le h^7/1024$ has two rows or two columns
that, after multiplication by a unit scalar, differ by at most $h$ in
every coordinate.

Dephase $E$ to rows $(1,1,1)$, $(x_0,x_1,x_2)$,
$(y_0,y_1,y_2)$. If all $x_j$ are within $h$ of $x_0$, the claim holds.
Otherwise choose $j,k$ with $|x_j-x_k|>h$ and interpolate
$y_j=\alpha+\beta x_j$, $y_k=\alpha+\beta x_k$.
Then $|\alpha|,|\beta|\le2/h$, and at the remaining index $l$,
\[
 |y_l-\alpha-\beta x_l|\le\eta:=|\det E|/h\le h/4.
\]
If $|\alpha|\le h/4$ or $|\beta|\le h/4$, normalizing the other
coefficient to modulus one makes the third row within
$2(h/4)+\eta\le3h/4$ of a multiple of the second or first row.

Otherwise put $L(x)=|\alpha+\beta x|^2-1$ on the unit circle.
The polynomial $xL(x)$ has roots $x_j,x_k$, so
\[
 |\alpha\beta|\,|x_l-x_j|\,|x_l-x_k|
 =|L(x_l)|\le3\eta.
\]
One distance $s=|x_l-x_j|$ or $|x_l-x_k|$ therefore satisfies
$s^2\le48|\det E|/h^3\le h^4/16$.
The corresponding $y$ distance is at most
$(2/h)s+\eta\le3h/4$. Those two columns are within $h$ entrywise.
This proves the phase-block estimate, including zero determinants.

Take $h=2^{-52}$. A corner with $|\det E|\le2^{-374}$ meets its
hypothesis. For a nearly proportional row pair, phase the first row
to ones and normalize the second at a selected column. The other two
relative phases have absolute value at most $4h=2^{-50}$.
The elementary three-phase routing argument in Section 5 of
\texttt{NEAR\_F6\_CLOSURE.md} then produces a dephased $F_3$ corner
with entry error at most $10\cdot2^{-50}\le2^{-46}$.
For a column pair apply the same argument to $H^T$; (17) is valid in
both orientations. Hence
\[
 M(H)\le5,\quad\min_c|\det H_c|\le2^{-374}
 \quad\Longrightarrow\quad d_F(H)<1/30.\tag{18}
\]

\subsection*{Step C: quantitative routing from nearly geometric rows}
Suppose three rows can be phased to $(1;x;y)$ with
$\max_j|y_j-\overline{x_j}|\le\delta=2^{-106}$.
Orthogonality gives $\sum x_j=0$ and $|\sum x_j^2|\le6\delta$.
If $e_k$ are the elementary symmetric functions of the six unit numbers
$x_j$, then $e_1=e_5=0$, $|e_2|\le3\delta$, and
$e_4=e_6\overline{e_2}$. Their root polynomial shows
\[
 |(x_j^3-U)(x_j^3-V)|\le6\delta,\tag{19}
\]
where $U,V$ are the roots of $T^2-e_3T+e_6$.
Put $\eta=\sqrt{6\delta}$. Assign each $x_j^3$ to one root within
$\eta$. Choosing a unit anchor $a$ or $b$ from each nonempty class
shows that every $x_j$ in that class is within $\eta$ of one of
$a,a\omega,a\omega^2$, or respectively $b,b\omega,b\omega^2$.
Here $\omega=e^{2\pi\ii/3}$; we used the elementary unit-circle bound
$\min_k|u-\omega^kv|\le|u^3-v^3|/2$.

If both classes occur, let $C_A,C_B$ be their count sums in
$\mathbb Z[\omega]$. The first two moment bounds give
\[
 |aC_A+bC_B|\le6\eta,\qquad
 |a^{-2}C_A+b^{-2}C_B|\le12\eta+6\delta.
\]
If $|a^3-b^3|>48\eta$, solving these two equations in the unknown
count sums gives $|C_A|,|C_B|<1/2$, since $\delta\le\eta$.
A nonzero Eisenstein integer has modulus at least one. Each count sum
therefore vanishes, so each nonempty class contains a complete
equilateral triple.

If $|a^3-b^3|\le48\eta$, rotate the second anchor by a cube root
to put it within $24\eta$ of the first. All six points are then within
$25\eta$ of one equilateral triple. Its total count sum has modulus
at most $150\eta<1$, so is zero; there are exactly two points of each
type. The same conclusion holds if only one class was nonempty.
Thus in all cases we can select three columns with
$x$ coordinates within $25\eta$ of an equilateral triple.

Using $y\approx\overline{x}$, then dephasing the selected block, gives
an $F_3$ corner with entry error at most $52\eta<2^{-46}$:
the last inequality is $52^2\cdot6<128^2$ after squaring.
Equation (17) proves
\[
 M(H)\le5,\quad\text{some geometric row-triple defect}\le2^{-106}
 \quad\Longrightarrow\quad d_F(H)<1/30.\tag{20}
\]
Any of the three rows may be the middle member; rephasing the columns
and reordering the rows reduces to the displayed form.

\subsection*{Step D: a positive minimum on the remaining strip domain}
Fix three rows and dephase them to $(1;x;y)$, with $x_0=y_0=1$.
They contain ten independent unit phases. Define the three defects
\[
 G_1=\sum_{j=1}^5|x_jy_j-1|^2,\quad
 G_2=\sum_{j=1}^5|y_jx_j^{-2}-1|^2,\quad
 G_3=\sum_{j=1}^5|x_jy_j^{-2}-1|^2.\tag{21}
\]
If any $G_k\le2^{-212}$, the relevant geometric-row error is at
most $2^{-106}$ and (20) applies. By (18), (20), a matrix with
$M(H)\le5$ outside the $F_6$ tube supplies a strip satisfying
\[
 |\det E_C|^2\ge2^{-748}\quad\text{for all twenty }C,
 \qquad G_k\ge2^{-212}\quad(k=1,2,3).\tag{22}
\]

On the compact set of orthogonal phase strips satisfying (22), the
function $w=\max_C|\mathcal D_{E_C}|$ has no zero.
Indeed, at a zero the critical-strip classification applies: all
minors are invertible and no geometric progression is possible.
The two exceptional rational endpoints are excluded by the strip
arguments in \texttt{CRITICAL\_P7\_EXCLUSION.md} and
\texttt{CRITICAL\_EQUI\_CORRELATION\_EXCLUSION.md}.
Every $p_C$ would consequently belong to $\{0,5,p_*\}$, contradicting
(12). These results apply to three orthogonal phase rows without
assuming that a completion exists: their proofs use only the six real
row-orthogonality equations, the complementary three-column blocks,
and their moment identities. No fourth row or completion is used.

For completeness we make the positive minimum explicit. Use the
$2^{10}$ Cayley charts
\[
 z=\sigma\frac{(1+\ii t)^2}{1+t^2},\qquad
 -1\le t\le1,\quad\sigma\in\{-1,1\},
\]
for the ten phases, and introduce $w\in[0,48]$.
Clear positive denominators in the six real orthogonality equations,
the forty inequalities $w\ge\pm\mathcal D_{E_C}$, the twenty minor
conditions and the three conditions (21)--(22). Together with the
twenty bounds on $t$ and two bounds on $w$, this gives 91 integer
polynomial constraints in 11 real variables.

Their degrees are at most 40 and their coefficient heights at most
$2^{800}$. Here are explicit sufficient bounds, using the sum of
absolute real and imaginary coefficients as a polynomial norm.
\begin{itemize}
\item Orthogonality: a common denominator uses at most ten variables
once, giving degree at most 20 and norm at most $6\cdot4^{10}<2^{23}$.
\item Signed invariants: for raw exponent $e\in\{-2,-1,0,1,2\}$,
denominator $(1+t^2)^2$ leaves numerator
$(1+\ii\operatorname{sgn}(e)t)^{2|e|}(1+t^2)^{2-|e|}$,
of degree four and norm at most $2^{2+|e|}$.
The 48-term census gives degree at most 37 including $w$ and norm
at most $2^{18}(36\cdot2^8+12\cdot2^{12})+2^{18}<2^{34}$.
\item Minors: a determinant uses at most six free entries. Its complex
numerator has degree at most 12 and norm at most $6\cdot4^6$.
Writing it $U+\ii V$ over the positive denominator $Q$, the condition
is $2^{748}(U^2+V^2)-Q^2\ge0$, of degree at most 24 and norm less
than $2^{779}$.
\item Geometry: use $Q=\prod_{v=1}^{10}(1+t_v^2)^2$.
Each term of $G_k$ is $2-m-\overline m$ with every phase exponent
bounded by two. Its cleared numerator has degree at most 40;
the constraint $2^{212}QG_k-Q\ge0$ has norm at most
$2^{212}\cdot5(2\cdot2^{20}+2\cdot2^{40})+2^{20}<2^{257}$.
\end{itemize}
All charts are compact, and $w$ is positive on each nonempty chart
domain, by the preceding strip argument.

Theorem 1 of Jeronimo--Perrucci--Tsigaridas,\footnote{
\href{https://arxiv.org/pdf/1112.0544}{On the minimum of a polynomial
function on a basic closed semialgebraic set and applications}, Theorem 1.}
applied with $(n,d,s,H)=(11,40,91,2^{800})$, bounds the positive minimum
of $w$ below by
\[
 \bigl(2^{4-11/2}\,2^{800}\,40^{11}\bigr)^{-11\cdot2^{11}\cdot40^{11}}
 >2^{-K},\qquad K=866\cdot11\cdot80^{11},\tag{23}
\]
because $2n+2s<2^{800}$ and $40^{11}<2^{66}$.
Take the connected component containing a global minimizer to meet
the theorem's compact-component hypothesis.

Finally, on a physical strip $\mathcal D_{E_C}^2=R_{E_{C^c}}$.
Thus its maximum squared is at most $M(H)$.
If $M(H)<2^{-2K}$, an outside-tube matrix would enter (22) and
contradict (23). This proves (14).

\paragraph{Interpretation.}
The global implication is now proved at the explicit cutoff $m_0$.
The exponent is about $1.64\cdot10^{25}$, compared with about
$7.50\cdot10^{52}$ in the earlier general confinement bound. Both
cutoffs are impractically small. The effective-minimum theorem replaces
the unproved stability estimate with a valid, extremely conservative
one; it does not show that the estimate remains valid up to 5.
The strip domain deliberately forgets the full carrier's distance from
$F_6$; a sharper minimum on that domain alone need not reach 5.
Reaching 5 requires larger routing neighbourhoods or retaining further
full-carrier information, for example constraints from several row strips
together with the distance condition.

\paragraph{Inputs and verification.}
The Fourier result uses only the exact census and its twelve explicit
Fourier equivalences, not MC1, the classification paper, a numerical
smoothness claim, or a no-four-MUBs certificate.
The tangent result uses the exact linearized Hadamard rank and signed
corner derivatives. The earlier compactness and imported
\texttt{EffectiveConfinement} statements use MC1 and its stated dependencies;
the new strip-stability theorem (14) has the separate inputs listed above.
Run \texttt{python3 -B simplerproof/averaging\_lemma/verify.py} from
the repository root. It recomputes the Fourier census, extremal census,
tangent forms and rational constants, and writes a receipt explicitly
marked as a partial result. The analytic inequalities in this note are
written proofs, not proof-assistant-checked statements.
For the follow-up stability theorem, run \texttt{verify\_stability.py}
in the same folder. It replays the critical-strip algebraic inputs and
the near-$F_3$ factorization/grid inputs, checks the new rational constants,
and audits the effective-minimum parameters. The routing arguments and
the application of the published theorem are the written proof above.
\end{document}
